\section{Introduction}
\label{sec:intro}

Fault localization tells a developer where a failing program is wrong. In everyday practice the trigger is a failing
test (run locally, in a pre-commit hook or on continuous integration) and, increasingly, the next step is to paste the
failure into an AI assistant: most developers now use or plan to use such assistants~\cite{stackoverflow2025,jetbrains2025},
debugging is among the tasks they bring to them~\cite{xiao2024devgpt}, and error information measurably helps LLMs
localize faults~\cite{wu2023llmfl}. With test-driven development~\cite{beck2003tdd}, those tests are
written first and state what the code is meant to do. Linters and type checkers do not help at this point when the fault
is in the logic: the code is well formed and well typed; it just computes the wrong value.

Research tools take a more elaborate route. Spectrum-based fault localization
(SBFL)~\cite{jones2002,abreu2009,wong2016survey} ranks statements by how often failing and passing tests execute them,
but it collapses on \emph{semantic} faults, where both kinds of test run the same lines. SemLoc~\cite{yang2026semloc}
targets exactly these faults: it asks an LLM for executable semantic constraints, instruments the program, runs the tests,
ranks lines by a constraint-violation spectrum, and patches suspicious lines to verify them. Its promise is
\emph{structured grounding}: the model proposes, execution decides. It reports \RqOneSemlocClaudeTZeroEightAuthorAllAccOne\%
Acc@1 on its SemFault-250 benchmark, far above SBFL.

\textbf{This paper asks whether that accuracy is grounded or guessed, and whether the pipeline beats the everyday
practice.} Answering requires more than rerunning the tool: one must check how results are scored, whether the
benchmark's labels point at the faults, how much each signal contributes, and whether conclusions survive real bugs,
renamed code and other models. We therefore ask four research questions.
\textbf{RQ1 (reproduction):} Do SemLoc's published results reproduce, and do they still hold under fair scoring, without
the guess bonus that the paper does not describe, and on labels checked by execution?
\textbf{RQ2 (comparison):} Is a single prompt with the failing tests at least as accurate as SemLoc, and at what cost?
\textbf{RQ3 (why):} Which signal drives the difference (failing tests or coverage), and do more elaborate additions help?
\textbf{RQ4 (robustness):} Do the conclusions hold (a)~on real bugs (BugsInPy~\cite{widyasari2020}), (b)~on identifier-renamed
programs, (c)~across five models and (d)~on all three label sets? Every LLM configuration runs three times, with program- and template-cluster bootstrap confidence intervals.

\noindent\textbf{Contributions.}
\begin{itemize}
\item \textbf{C1 (RQ1), guessed by design and mislabeled:} SemLoc's top-1 line is, by an undocumented choice, the LLM's own
guess; its execution evidence acts below the first line; its counterfactual step (\CostCfSharePct\% of LLM calls) never
changes a ranking. We also report a released scoring script that credits delta debugging with \RqOneDdAuthorAllAccOne\%
where the paper reports zero, two bugs in the released BugsInPy runner, and a re-execution that falls short of the logs.
\LabaudAllMoved{} of \LabaudAllN{} labels move under review; we release an execution-repaired label set (tests decide) and
a reviewed one (an LLM pre-review, then team review).
\item \textbf{C2 (RQ2), the simple prompt wins:} one prompt with failing tests beats SemLoc at Acc@1 by
\RqTwoRepairedAllLlmTestsDAccOne{} [\RqTwoRepairedAllLlmTestsDAccOneLo, \RqTwoRepairedAllLlmTestsDAccOneHi] and at Acc@3
by \RqTwoRepairedAllLlmTestsDAccThree{} points on repaired labels, at about a seventh of the cost.
\item \textbf{C3 (RQ3), why, and nothing more elaborate closes the gap:} failing tests carry the signal and per-line
coverage does not; CovGate (coverage added to SemLoc) helps SemLoc only slightly, and propose-and-verify (test-passing
patches as arbiter), voting and static analysis do not improve on the simple prompt, because small patches pass the
tests at many non-faulty candidate lines.
\item \textbf{C4 (RQ4), robustness:} the simple prompt still beats SemLoc on real BugsInPy bugs (Acc@3
\BipAllLlmTestsAccThreeMean\% vs.\ \BipAllSemlocPriorAccThreeMean\%), barely moves on identifier-renamed programs (Acc@3
\RenRepairedLlmTestsDAccThree{} points), and does not lose to SemLoc with any of five models.
\end{itemize}
Alongside, we provide a seven-item evaluation checklist and a fully replayable artifact, including every reviewer's
answers. Code, data and LLM logs: \url{https://github.com/doanan293/grounded-or-guessed}.

Section~\ref{sec:related} reviews related work, Section~\ref{sec:background} introduces SBFL and SemLoc,
Section~\ref{sec:approach} describes the audit, the label sets, the grounded-vs-guessed decomposition and the methods, Section~\ref{sec:eval} answers the four
research questions, Section~\ref{sec:discussion} discusses the findings and threats, and Section~\ref{sec:conclusion} concludes.
